% Generated draft: requires mathematical and source review.
\documentclass[11pt]{article}
\usepackage[margin=25mm]{geometry}
\usepackage{fontspec}
\setmainfont{Latin Modern Roman}
\usepackage{amsmath,amssymb,mathtools,unicode-math}
\setmathfont{Latin Modern Math}
\usepackage{enumitem,needspace,xurl,hyperref}
\setlength{\parindent}{0pt}
\setlength{\parskip}{5pt}
\setlength{\emergencystretch}{3em}
\newcommand{\N}{\mathbb{N}}
\newcommand{\Z}{\mathbb{Z}}
\newcommand{\Q}{\mathbb{Q}}
\newcommand{\R}{\mathbb{R}}
\newcommand{\C}{\mathbb{C}}
\newcommand{\F}{\mathbb{F}}
\newcommand{\A}{\mathbb{A}}
\newcommand{\PP}{\mathbb{P}}
\newcommand{\E}{\mathbb{E}}
\newcommand{\eps}{\varepsilon}
\newcommand{\dd}{\,\mathrm{d}}
\begin{document}
\section*{UnsolvedMath ID: 701592}
\textbf{Source problem code:} catalogue entry
\subsection{Definitions and mathematical statement}
For elements of a group, put \([x,g]=x^{-1}g^{-1}xg\), and define the left-normed iterated Engel commutators by
\[
[x,{}_1g]=[x,g],\qquad [x,{}_{k+1}g]=[[x,{}_kg],g]\quad(k\geq1).
\]
Following the supplied source, a profinite group has rank at most \(r\) if every subgroup can be topologically generated by at most \(r\) elements; it has finite rank if such an integer \(r\geq0\) exists. For a subgroup that is not closed, topological generation may be interpreted in its induced topology: generators \(h_1,\ldots,h_s\in H\), with \(s\leq r\), satisfy \(\overline{\langle h_1,\ldots,h_s\rangle}^{\,H}=H\). The source does not explicitly distinguish this reading from a definition quantified only over closed subgroups; this convention requires review.

Let \(G\) be a profinite group satisfying
\[
\forall g\in G\;\exists E_g\leq G\text{ closed of finite rank}\;\forall x\in G\;\exists n(x,g)\geq1\;\forall k\geq n(x,g),
\qquad [x,{}_kg]\in E_g.
\]
The finite rank bound for \(E_g\) may depend on \(g\); the threshold may depend on both \(x\) and \(g\). No uniform bound is assumed.

Must there exist a normal subgroup \(N\trianglelefteq G\) of finite rank such that \(G/N\) is locally nilpotent? Here local nilpotence means that every finitely generated abstract subgroup of \(G/N\) is nilpotent, with no uniform bound on its nilpotency class.

The source does not explicitly require \(N\) to be closed. Consequently, whether the intended conclusion includes closedness of \(N\), and hence a profinite quotient with its quotient topology, remains unresolved. If nonclosed \(N\) is allowed, its finite rank must also be interpreted using its induced topology, since the supplied definition is stated for profinite groups.
\subsection{Short English statement}
Suppose each element \(g\) of a profinite group has a closed subgroup of finite rank containing all sufficiently long Engel commutators with \(g\), with the required length allowed to depend on the starting element. Must the group have a normal subgroup of finite rank whose quotient is locally nilpotent? The source leaves the closedness of this normal subgroup and the precise subgroup convention in the rank definition unstated.
\subsection{Sources}
\begin{itemize}
\item Local UnsolvedMath registry, supplied complete source record: permanent ID \(701592\), exact source code \texttt{catalogue entry}, title \emph{Kourovka Notebook Problem 19.44}; upstream revision \texttt{d371ee917ac9647fbb06cd6fc529deec6f3dda82}. The supplied registry labels the problem open; this label is not an independently verified mathematical assessment.
\item \emph{The Kourovka Notebook (Unsolved Problems in Group Theory)}, Problem 19.44, proposed by E. I. Khukhro in 2018. The supplied background identifies the 19th issue (2018) and the 21st issue edition (2026), with identifier \texttt{arXiv:1401.0300v48}. The registry describes the collection as edited by E. I. Khukhro and V. D. Mazurov. These bibliographic details and the problem's attribution are supplied source information and have not been independently verified.
\end{itemize}
\end{document}
